\documentclass{gmto}

\usepackage{amsmath}
\usepackage{amssymb}
\usepackage{booktabs}
\usepackage[section,above]{placeins}
\numberwithin{table}{section}
\numberwithin{equation}{section}

\graphicspath{{../figures/}}
\tolerance=2000
\emergencystretch=3em

\newcommand{\code}[1]{\texttt{#1}}
\newcommand{\mas}{\,\mathrm{mas}}
\newcommand{\arcsec}{^{\prime\prime}}
\newcommand{\vect}[1]{\boldsymbol{#1}}

\DocID{GMT-XXX-\#\#\#\#}
\DocVersion{Ver. X0}
\DocStatus{Draft}

\title{Mount Kinematics and Pointing Errors in Tracking Simulations}
\subtitle{GMT Technical Note --- \code{trajectory-servos-mount}}
\author{Claude Opus 5.5 (Anthropic)}
\date{10/20/2026}

\begin{document}

\maketitle

\clearpage

\section*{Signatures}
\vspace{1cm}
\subsection*{Author}
\vspace{1.5cm}
\begin{tabu} to \linewidth {X[3,l]X[1,l]}
  \rule{\linewidth}{.1pt} & \rule{\linewidth}{.1pt} \\
  Claude Opus 5.5, AI model (Anthropic) & Date
\end{tabu}
\vspace{1.5cm}
\subsection*{Approvers}
\vspace{1.5cm}
\begin{tabu} to \linewidth {X[3,l]X[1,l]}
  \rule{\linewidth}{.1pt} & \rule{\linewidth}{.1pt} \\
  Rod Conan               & Date                    \\[1cm]
\end{tabu}

\clearpage

\section*{Revision Log}

\begin{revisions}
  X0 & 10/20/2026 & All & None & & Claude Opus 5.5 \\
\end{revisions}

\clearpage

\tableofcontents
\listoffigures
\listoftables

\clearpage

%======================================================================
\section{Purpose}
\label{sec:purpose}

This note describes how the telescope motion along a sidereal tracking trajectory is modeled in the
\code{trajectory-servos-mount} demonstration of the GMT integrated modeling framework
(\code{dos-actors}, \code{demos/src/bin/trajectory-servos-mount}).
It explains, and gives the equations implemented for:
\begin{itemize}
\item the generation of the tracking trajectory (Section~\ref{sec:trajectory}),
\item the mount drives feed-forward of the FEM rigid body springs (Section~\ref{sec:feedforward}),
\item the kinematic compensation of the rigid telescope motion in the FEM outputs (Section~\ref{sec:kinematics}),
\item the rotation of the pupil as seen by an instrument on the Gregorian Instrument Rotator (GIR)
  (Section~\ref{sec:pupil}),
\item the telescope pointing and field rotation errors fed to the optical model (Section~\ref{sec:pointing}),
\end{itemize}
and it reports the verifications that were performed on each of them (Section~\ref{sec:verification}).

\subsection{Introduction}
\label{sec:introduction}

The integrated model couples the GMT structural finite element model (FEM), the mount control
system, the M1 and M2 control systems (\code{GmtServoMechanisms}), a linear optical model (LOM) and,
optionally, the \code{crseo} ray tracing optical model.
The mount is driven along the 5 minutes ATP default tracking trajectory and the model is sampled
at the mount control rate, 1\,kHz for the default \code{MOUNT\_FDR\_1kHz} mount model.
The FEM is the 2025 model
\code{20251110\_1617\_zen\_30\_M1\_202110\_FSM\_202305\_}\linebreak\code{Mount\_202305\_pier\_202411\_M1Fans}
computed at 30$^\circ$ zenith angle.

The FEM is a linear model: a rotation of the mount by an angle $\theta$ appears in the FEM outputs as
the linearized rigid displacement of every node of the telescope, superimposed on the structure
deformations. Along a 300\,s trajectory the azimuth axis rotates by 1.7$^\circ$, which moves
the M1 outer segments by about 0.25\,m, far beyond the range of validity of a linear optical model.
Three things are therefore required:
\begin{enumerate}
\item the mount drives must be able to hold the FEM away from its reference position, whose
  ``rigid body'' modes are in fact very soft springs (Section~\ref{sec:feedforward});
\item the rigid motion must be removed from the FEM outputs, leaving the deformations only
  (Section~\ref{sec:kinematics});
\item the part of the rigid motion that \emph{does} matter optically, i.e.\ the difference between where
  the telescope points and where it should point, and the rotation of the instrument relative to
  the pupil and to the sky, must be reintroduced explicitly in the optical model
  (Sections~\ref{sec:pupil} and~\ref{sec:pointing}).
\end{enumerate}

\subsection{Model overview}
\label{sec:overview}

Table~\ref{tab:clients} lists the clients of the model involved in this note and
Table~\ref{tab:signals} the signals exchanged between them.

\begin{table}[htbp]
  \centering
  \caption{Clients of the \code{trajectory-servos-mount} model.}
  \label{tab:clients}
  \begin{tabularx}{\linewidth}{lX}
    \toprule
    Client (source file)                                      & Role \\
    \midrule
    \code{Trajectory} (\code{main.rs})                        & ATP trajectory, mount setpoints and field rotation \\
    \code{SpringFeedForward} (\code{feedforward.rs})          & mount drive torques holding the FEM at the setpoint angles \\
    \code{GmtServoMechanisms}                                 & FEM with kinematic compensation, mount, M1 and M2 control systems \\
    \code{LinearOpticalModel}                                 & segment tip-tilt from the M1 and M2 rigid body motions \\
    \code{GirRotation} (\code{gir\_rotation.rs})              & M1 and M2 rigid body motions rotated by the GIR angle \\
    \code{Pointing} (\code{pointing.rs})                      & telescope pointing and rotation errors from the mount tracking error \\
    \code{OpticalModel<Camera<1>>} & ray tracing optical model and imaging camera \\
    \bottomrule
  \end{tabularx}
\end{table}

\begin{table}[htbp]
  \centering
  \caption{Signals of the model relevant to this note. Rates are given for the 1\,kHz mount model.}
  \label{tab:signals}
  \begin{tabularx}{\linewidth}{llX}
    \toprule
    Signal                          & Rate     & Content \\
    \midrule
    \code{MountSetPoint}            & 1\,kHz   & azimuth, elevation and GIR setpoints [rad] \\
    \code{FieldRotation}            & 1\,kHz   & change of the parallactic angle [rad] \\
    \code{MountTorquesFeedForward}  & 1\,kHz   & 8 azimuth, 4 elevation and 4 GIR drive torques [N\,m] \\
    \code{MountEncoders}            & 1\,kHz   & all the mount encoders, GIR encoders at indices 10--13 \\
    \code{AverageMountEncoders}     & 1\,kHz   & averaged azimuth, elevation and GIR encoders [rad] \\
    \code{M1/M2RigidBodyMotions}    & 1\,kHz / 50\,Hz & $7\times 6$ segment rigid body motions in the segment local frames \\
    \code{PointingError}            & 50\,Hz   & telescope optical axis relative to the source, $(x,y)$ in the OSS [rad] \\
    \code{PointingRotationError}    & 50\,Hz   & instrument rotation about the optical axis relative to the sky [rad] \\
    \code{TipTilt}                  & 50\,Hz   & wavefront tip-tilt, ordered $[y,x]$ [rad] \\
    \code{Frame<Host>}              & 50\,Hz   & $256\times 256$ camera frame \\
    \bottomrule
  \end{tabularx}
\end{table}

%======================================================================
\section{Definitions, Acronyms, and Reference Documents}
\label{sec:definitions-acronyms-references}

\subsection{Definitions}
\label{sec:definitions}

\begin{definitions}
  OSS & Optical Support Structure coordinate system: $z$ along the optical axis, $x$ along the elevation axis \\
  Tracking error & mount encoders minus mount setpoint \\
  Field rotation & change of the parallactic angle along the trajectory, followed by the GIR \\
  Kinematic gain & static gain $K$ from the averaged mount encoder angles to the FEM outputs \\
  Absolute pointing & camera pointing fixed with respect to the optical axis, the image moving with the source direction
\end{definitions}

\subsection{Acronyms}
\label{sec:acronyms}

\begin{acronyms}
  ATP  & Acquisition, Tracking and Pointing (TBC) \\
  CEO  & CUDA Engined Optics \\
  DC   & Direct Current (static, zero frequency) \\
  DFT  & Discrete Fourier Transform \\
  FEM  & Finite Element Model \\
  FOV  & Field Of View \\
  GIR  & Gregorian Instrument Rotator \\
  GMT  & Giant Magellan Telescope \\
  GMTO & Giant Magellan Telescope Corporation\\
  ICRS & International Celestial Reference System \\
  LOM  & Linear Optical Model \\
  LST  & Local Sidereal Time \\
  OSS  & Optical Support Structure \\
  RBM  & Rigid Body Motion \\
  TETE & True Equator True Equinox
\end{acronyms}

\subsection{Applicable Documents}
\label{sec:aplicables}

\begin{documents}[Applicable Documents]
  TBD && \\
\end{documents}

\subsection{Reference Documents}
\label{sec:references}

\begin{documents}[Referenced Documents]
  GMT-REF-00362 & GMT Project Acronyms and Glossary & \url{https://bit.ly/3l4reUC} \\
  RD1 & \code{trajectory-servos-mount} demo: README, sources and scripts (\code{demos/src/bin/}) & \code{dos-actors} repository \\
  RD2 & \code{gmt\_dos-clients\_fem} kinematic compensation (\code{kinematics.rs}) & \code{dos-actors} repository \\
  RD3 & \code{gmt\_dos-clients\_crseo} pointing errors (\code{optical\_model.rs}) & \code{dos-actors} repository \\
\end{documents}

%======================================================================
\section{Tracking trajectory}
\label{sec:trajectory}

\subsection{Computation with astropy}
\label{sec:astropy}

The trajectory \code{atp\_trajectory.csv} is computed by \code{scripts/atp\_trajectory.py}
(\code{make trajectory}) with astropy for the ATP default pointing:
\begin{itemize}
\item site: Las Campanas, latitude $\phi = -29.049^\circ$, longitude $-70.682^\circ$, height 2514\,m;
\item epoch: $t_0 = $ 2018-01-01T04:00:00 UTC;
\item initial pointing: altitude 45$^\circ$, azimuth 0$^\circ$, rotator 0$^\circ$.
\end{itemize}
The sky target is fixed by converting the initial altitude and azimuth at $t_0$ to ICRS coordinates,
without atmospheric refraction ($\alpha = 89.8175^\circ$, $\delta = 15.9532^\circ$).
The target is then transformed back to the local horizontal frame (\code{AltAz}) at 301 epochs,
$t_0 + k$\,s, $k = 0,\dots,300$, giving the elevation and the azimuth, the latter unwrapped to
avoid the $\pm 180^\circ$ discontinuity.

The rotator angle is the change of the parallactic angle
\begin{equation}
  q = \operatorname{atan2}\!\left(\sin H,\; \tan\phi\cos\delta - \sin\delta\cos H\right),
  \label{eq:parallactic}
\end{equation}
where the hour angle $H$ is the apparent local sidereal time minus the target right ascension
in the TETE frame (equinox based apparent place, consistent with the apparent LST), and
$\delta$ the TETE declination. $q$ is unwrapped and the rotator angle is $q(t) - q(t_0)$.

The file holds 4 columns: time [s], azimuth, elevation and rotator [deg].
Over 300\,s the azimuth changes by $-1.704^\circ$ ($-6134\arcsec$, about $-20.4\arcsec$/s),
the elevation by $-58.7\arcsec$ and the rotator by $-1.549^\circ$ ($-5577\arcsec$, about
$-18.6\arcsec$/s).

\subsection{Mount setpoints and field rotation}
\label{sec:setpoints}

The FEM is linearized about its reference configuration, so the mount setpoints are the
offsets of the trajectory from its first sample. The \code{Trajectory} client reads the file,
converts the angles to radians, subtracts the first sample, and linearly interpolates the 1\,Hz
samples to the mount sampling frequency $f_s$: at step $n$, with $i = \lfloor n/f_s\rfloor$ and
$a = (n \bmod f_s)/f_s$,
\begin{equation}
  \vect{s}(n) = \vect{s}_i + a\,(\vect{s}_{i+1} - \vect{s}_i),
\end{equation}
where $\vect{s}_i = [\mathrm{az}_i, \mathrm{el}_i, q_i]$.
It writes two outputs:
\begin{itemize}
\item \code{MountSetPoint}: $[\mathrm{az}, \mathrm{el}, q]$, or $[\mathrm{az}, \mathrm{el}, 0]$
  when \code{GIR\_TRACKING=0} (the GIR is at rest);
\item \code{FieldRotation}: $q$, in both cases, since the field rotates whether or not the GIR
  follows it.
\end{itemize}

%======================================================================
\section{Mount drives feed-forward}
\label{sec:feedforward}

\subsection{Why it is needed}
\label{sec:ff-why}

In the FEM, the azimuth, elevation and GIR ``rigid body'' modes are not free: they have small but
non-zero eigen frequencies, between 0.07\,Hz and 0.1\,Hz. Each mount axis therefore behaves as a
very soft spring, and holding it at an angle $\theta$ away from the FEM reference requires a
static drive torque proportional to $\theta$. In the mount controller this torque has to be built
up by the integral action of the feedback loop. Along the trajectory the azimuth angle grows
linearly, so does the required torque, and the azimuth drives reach their 800\,kN\,m saturation
at about 242\,s: the mount then stops tracking. The spring is an artifact of the FEM, a real mount
does not need a torque to hold an angle, so it is cancelled by a feed-forward rather than by
retuning the controller.

\subsection{How it works}
\label{sec:ff-how}

\code{SpringFeedForward::from\_fem} derives, from the FEM, the torque per drive and per radian of
each axis. The static (DC) gain from the drive torques $\vect{\tau}$ (8 azimuth, 4 elevation and 4 GIR
drives) to any FEM output $\vect{y}$ is computed from all the FEM modes,
\begin{equation}
  \vect{y} = \Phi\,\Omega^{-2}\,B\,\vect{\tau} = G\,\vect{\tau},
  \label{eq:dcgain}
\end{equation}
with $B$ the inputs to modal forces matrix (\code{in2modes}) of the drive torques, $\Omega^2$ the
diagonal matrix of the squared modal angular frequencies $\omega_i^2$ and $\Phi$ the modes to
outputs matrix (\code{modes2out}).
Applying Eq.~(\ref{eq:dcgain}) to the encoders of each axis and averaging over the encoders of the
axis gives the gain $g_{k,j}$ from drive $j$ to the averaged encoder angle of axis
$k \in \{\mathrm{az}, \mathrm{el}, \mathrm{GIR}\}$.
If the same torque $\tau_k$ is applied on the $n_k$ drives $\mathcal{D}_k$ of axis $k$,
the axis moves by $\theta_k = \tau_k\sum_{j\in\mathcal{D}_k} g_{k,j}$, so the torque per drive that
holds the axis at the setpoint $\theta_k$ is
\begin{equation}
  \tau_k = \kappa_k\,\theta_k, \qquad
  \kappa_k = \Big(\sum_{j\in\mathcal{D}_k} g_{k,j}\Big)^{-1}.
  \label{eq:stiffness}
\end{equation}
At each step the client reads \code{MountSetPoint} and writes \code{MountTorquesFeedForward},
the 16 torques $[\kappa_{\mathrm{az}}\theta_{\mathrm{az}}\times 8,\;
\kappa_{\mathrm{el}}\theta_{\mathrm{el}}\times 4,\; \kappa_{\mathrm{GIR}}\theta_{\mathrm{GIR}}\times 4]$.
The mount controller (\code{gmt\_dos-clients\_mount}) adds them to the drive torques
\emph{after} the drives saturation, so the saturation only applies to the feedback torques.
The cross-coupling between axes is neglected: only the diagonal terms $g_{k,j}$, $j\in\mathcal{D}_k$,
are used; the coupling is estimated in Section~\ref{sec:ver-ff} and found negligible. The stiffness $[\kappa_{\mathrm{az}}, \kappa_{\mathrm{el}}, \kappa_{\mathrm{GIR}}]$
is printed at start-up; for the 2025 FEM it is
$[4.17\times10^{6},\; 1.08\times10^{7},\; 1.00\times10^{5}]$\,N\,m/rad per drive. \code{MOUNT\_FEEDFORWARD=0}
turns the feed-forward off.

%======================================================================
\section{Kinematic compensation}
\label{sec:kinematics}

\subsection{Principle}
\label{sec:kin-principle}

As the mount moves, every FEM output $\vect{y}$ (segment rigid body motions, hardpoint forces,
M2 positioners, \ldots) includes the motion that follows the mount, i.e.\ the rigid telescope motion
plus the static deflection of the structure that comes with it. The kinematic compensation removes
from each output the part that is linearly related to the averaged mount encoder angles
$\vect{\theta} = [\bar\theta_{\mathrm{az}}, \bar\theta_{\mathrm{el}}, \bar\theta_{\mathrm{GIR}}]$:
\begin{equation}
  \vect{y}_c = \vect{y} - K\,\vect{\theta},
  \label{eq:kin}
\end{equation}
where $K$ is the $n_y\times 3$ kinematic gain.
The compensation is applied in the FEM client after each state update
(\code{DiscreteModalSolver::update}), so that every consumer of the FEM outputs, including the M1 and
M2 control systems, sees compensated outputs.

\subsection{Kinematic gain from the FEM DC gain}
\label{sec:kin-gain}

$K$ is derived from the same static gain as the feed-forward, Eq.~(\ref{eq:dcgain}), computed in
\code{DiscreteStateSpace::kinematic\_compensator}.
Let $H$ be the $n_y\times 16$ DC gain from the 16 drive torques to all the FEM outputs, and $E$ the
$3\times 16$ DC gain from the drive torques to the averaged azimuth, elevation and GIR encoders (the
row mean of $H$ over the encoders of each axis). A static torque $\vect{\tau}$ produces
$\vect{y} = H\vect{\tau}$ and $\vect{\theta} = E\vect{\tau}$. The kinematic gain is the
least-squares map from $\vect{\theta}$ to $\vect{y}$ over the torques:
\begin{equation}
  K = H E^{+}, \qquad E^{+} = E^{T}\left(E E^{T}\right)^{-1}.
  \label{eq:kgain}
\end{equation}
Each column of $K$ is the static response of the outputs per radian of one mount axis, as measured
by its encoders. Because $K$ is computed from all the FEM modes, it contains the exact linearized
rigid motion of the FEM (with its actual geometry and axes) and the quasi-static flexure that goes
with it, with no geometric model of the telescope.

The rows of $K$ of the mount encoders are set to zero: the mount control system needs the
uncompensated encoders. Other outputs can be left uncompensated
(\code{kinematic\_exclusion}), and a column of $K$ can be replaced by a calibrated gain
(\code{kinematic\_calibration}); neither is used in the demo, where
\code{GmtServoMechanisms::kinematic\_compensation} compensates all the outputs but the mount encoders.

\subsection{Compensation of the control feedback}
\label{sec:kin-feedback}

The outputs fed back to the control systems must be compensated too. The M1 hardpoint forces, in
particular, contain a kinematic component as the telescope rotates: the M1 actuators force loop would
offload it and make the M1 segments drift with the mount azimuth (Section~\ref{sec:ver-kin}).
Compensating all the FEM outputs removes the drift without closed-loop calibration.

\subsection{Consequence for the optics}
\label{sec:kin-optics}

After compensation the M1 and M2 rigid body motions only contain the structure deformations
relative to the mount encoders (Figure~\ref{fig:rbm}). The rigid telescope motion, including the
tracking error, has been removed with the encoders. That is the desired behavior for the large
trajectory motion, which the sky follows, but it also removes the tracking error, which
does move the image. The tracking error is therefore reintroduced as a pointing error in the optical
model (Section~\ref{sec:pointing}). \code{KINEMATIC\_COMPENSATION=0} runs the FEM without compensation.

%======================================================================
\section{Pupil rotation}
\label{sec:pupil}

\subsection{Principle}
\label{sec:pupil-principle}

The FEM outputs are given in the OSS, which is fixed to the telescope tube. An instrument mounted on
the GIR rotates with it about the optical axis by the GIR angle $\gamma$, and sees the pupil (the M1
and M2 segments) rotated with respect to its own frame. With the optical model feature
(\code{gmt\_dos-clients\_crseo}), \code{GirRotation} (\code{gir\_rotation.rs}) rotates the M1 and
M2 segments as rigid bodies about the OSS $z$ axis by $\gamma$, the mean of the 4 GIR encoders
(\code{MountEncoders} indices 10--13, uncompensated), before ray tracing.

\subsection{Implementation}
\label{sec:pupil-impl}

Consider a segment with vertex $\vect{r}$ and local frame $Q$ ($\vect{v}_{\mathrm{OSS}} = Q\,\vect{v}_{\mathrm{local}}$),
deformed by the compensated rigid body motion $(\vect{t}_d, R_d)$ from the FEM, given in its local frame.
The deformed segment vertex is at $\vect{r} + Q\vect{t}_d$ and its orientation is $QR_d$ in the OSS.
Rotating the deformed mirror by $R_z(\gamma)$ and expressing the result in the nominal local frame gives
the rigid body motion fed to the optical model:
\begin{align}
  \vect{t} &= Q^{T}\left(R_z(\gamma) - I\right)\vect{r} + M\,\vect{t}_d, \label{eq:gir-t}\\
  R &= M R_d, \qquad M = Q^{T} R_z(\gamma)\, Q, \label{eq:gir-r}
\end{align}
$M$ being the mirror rotation in the segment local frame. The rotations are given as the $(x,y,z)$ Euler
angles of the ray tracing coordinate systems, $R = R_z(c)R_y(b)R_x(a)$, from which
\begin{equation}
  b = -\arcsin R_{20},\quad a = \operatorname{atan2}(R_{21}, R_{22}),\quad
  c = \operatorname{atan2}(R_{10}, R_{00}).
\end{equation}
The deformation thus rotates with the mirror, as it does seen from an instrument on the GIR.
A first implementation added the rotation of the undeformed mirror to the deformation,
$\vect{t} = Q^{T}(R_z(\gamma) - I)\vect{r} + \vect{t}_d$ and the sum of the Euler angles of $M$ and $R_d$:
the deformation was left in the orientation of the telescope tube, an error of the order of
$\gamma$ times the deformation. Over the 30\,s run of Section~\ref{sec:ver-e2e} ($\gamma \le 0.15^\circ$),
the exact composition changes the wavefront tip-tilt by $1.3\times10^{-3}\mas$ on average and at most
$6\times10^{-3}\mas$; at the end of the 300\,s trajectory ($\gamma = 1.55^\circ$) the error would have
been about 2.7\% of the deformation terms.
The segment geometry (Table~\ref{tab:geometry}) follows the geotrans conventions: the 6 outer
segments at a radial distance $\rho$ with clocking $o_i = o_0 - 60^\circ i$, the vertex height being
the mirror height plus the conic sag $z(\rho) = \rho^2/\big(R(1+\sqrt{1-(1+k)\rho^2/R^2})\big)$, and
the local frames $Q_i = R_z(o_i)R_x(\beta)$ for M1 and $Q_i = R_z(o_i)R_y(\pi)R_x(\beta)$ for M2.
The central segment (S7) frame is the identity for M1 and $R_z(\pi)R_y(\pi)$ for M2,
whose S7 frame is flipped with respect to the OSS.

\begin{table}[htbp]
  \centering
  \caption{Segment geometry used by \code{GirRotation}.}
  \label{tab:geometry}
  \begin{tabular}{lcc}
    \toprule
                                 & M1            & M2 \\
    \midrule
    outer segment distance $\rho$ [m] & 8.71     & 1.08774 \\
    clocking $o_0$ [deg]         & 0             & 180 \\
    vertex height [m]            & 3.9           & 3.9 + 20.26247614 \\
    radius of curvature $R$ [m]  & 36            & $-4.1639009$ \\
    conic constant $k$           & $-0.9982857$  & $-0.71692784$ \\
    outer segment tilt $\beta$ [deg] & 13.601685 & 14.777498 \\
    \bottomrule
  \end{tabular}
\end{table}

The rotated rigid body motions are fed to the optical model at 50\,Hz. Rotating the pupil rotates
the diffraction pattern of the segmented aperture (and any segment piston/tip-tilt pattern) on the
detector; it does not move the image centroid of a perfectly phased telescope, which is set by the
source direction (Section~\ref{sec:pointing}).

%======================================================================
\section{Telescope pointing errors}
\label{sec:pointing}

\subsection{Pointing error from the mount tracking error}
\label{sec:pointing-error}

The \code{Pointing} client (\code{pointing.rs}) reads \code{MountSetPoint},
\code{AverageMountEncoders} and \code{FieldRotation} at 1\,kHz.
The tracking errors are $\Delta\mathrm{az} = \bar\theta_{\mathrm{az}} - s_{\mathrm{az}}$ and
$\Delta\mathrm{el} = \bar\theta_{\mathrm{el}} - s_{\mathrm{el}}$.
In the OSS, with $\zeta$ the zenith angle, the azimuth axis is $(0, \sin\zeta, \cos\zeta)$ and the
elevation axis is $(-1, 0, 0)$. The telescope rotation vector is
\begin{equation}
  \vect{w} = \Delta\mathrm{az}\,(0, \sin\zeta, \cos\zeta) + \Delta\mathrm{el}\,(-1, 0, 0),
\end{equation}
which moves the optical axis $\hat{\vect{z}}$ by
$\vect{w}\times\hat{\vect{z}} = (w_y, -w_x, 0)$, i.e.\ the pointing error in the OSS is
\begin{equation}
  \vect{p}_{\mathrm{OSS}} = \left(\Delta\mathrm{az}\cos(\mathrm{el}),\; \Delta\mathrm{el}\right).
  \label{eq:pe}
\end{equation}
The elevation is the telescope elevation, $\mathrm{el} = \mathrm{el}_0 + \bar\theta_{\mathrm{el}}$,
with $\mathrm{el}_0$ the initial elevation of the trajectory (45$^\circ$, read from the first row of
the trajectory file) and $\bar\theta_{\mathrm{el}}$ the averaged elevation encoder, the encoders being
offsets from the initial pointing. Along the ATP trajectory the elevation decreases by $58.7\arcsec$,
changing $\cos(\mathrm{el})$ by $3\times10^{-4}$ in relative terms.
As the mirrors (Section~\ref{sec:pupil}), the pointing error is rotated by the GIR angle
$\gamma = \bar\theta_{\mathrm{GIR}}$ into the instrument frame:
\begin{equation}
  \vect{p} = R(\gamma)\,\vect{p}_{\mathrm{OSS}}, \qquad
  R(\gamma) = \begin{pmatrix}\cos\gamma & -\sin\gamma\\ \sin\gamma & \cos\gamma\end{pmatrix}.
\end{equation}
$\vect{p}$ is the \code{PointingError} output.

\subsection{Field rotation error}
\label{sec:rotation-error}

The GIR compensates the field rotation $q$ (Section~\ref{sec:setpoints}). Its tracking error relative
to the sky,
\begin{equation}
  \Delta\gamma = \bar\theta_{\mathrm{GIR}} - q,
\end{equation}
is the \code{PointingRotationError} output. With the GIR tracking, $\Delta\gamma$ is the GIR tracking
error; with \code{GIR\_TRACKING=0}, the GIR is at rest and $\Delta\gamma = -q$: the field rotates on
the instrument. \code{POINTING\_ERROR=0} sets both outputs to zero, and
\code{POINTING\_ROTATION\_GAIN} scales $\Delta\gamma$ for verification only
(Section~\ref{sec:ver-e2e}).
Both outputs are decimated to 50\,Hz, the rate of the optical model.

\subsection{Optical model}
\label{sec:pointing-om}

The optical model images a J band source (1.215\,\textmu m) at an off-axis angle $\theta_0$
(\code{OFF\_AXIS}, default $0.25\arcsec$ along $x$) on a camera with a $\lambda/(5D)$ DFT sampling
binned $2\times 2$, i.e.\ a $3.93\mas$ pixel over a $256\times 256$ frame ($\pm 0.5\arcsec$ FOV).
The camera is in \emph{absolute pointing} along the optical axis
(\code{sensor\_pointing}$(-\theta_0, 0)$): the image sits at the source direction seen from the
instrument, $\theta_0$ from the frame center, and moves when that direction changes.
This requires \code{crseo-sys}~1.3.3 or later (CEO imaging fix, the $x$ axis of the absolute pointing
along the frame columns).

When the optical model reads \code{PointingError} or \code{PointingRotationError}, the source
directions are set to the directions seen from the instrument
(\code{Pointing::apply} in \code{gmt\_dos-clients\_crseo}): for a nominal source at
$(x, y) = \zeta_s(\cos a_s, \sin a_s)$,
\begin{equation}
  \begin{pmatrix} x' \\ y' \end{pmatrix} =
  R(-\Delta\gamma)\begin{pmatrix} x \\ y \end{pmatrix} - \vect{p},
  \label{eq:src}
\end{equation}
the nominal directions being kept unchanged. The source is ray traced along $(x', y')$, so the image
moves by $-\vect{p}$ and is rotated by $-\Delta\gamma$ about the optical axis. For the off-axis source,
the rotation moves the image by $\simeq -\Delta\gamma\,\theta_0$ perpendicular to its field
direction. The wavefront tip-tilt (\code{TipTilt}, ordered $[y, x]$) is computed relative to the
source direction, and thus contains the structure deformations only. The expected image motion is
\begin{equation}
  \delta\vect{c} = \big(R(-\Delta\gamma) - I\big)\begin{pmatrix}\theta_0\\0\end{pmatrix}
  - \vect{p} + \vect{t}_{\mathrm{TT}},
  \label{eq:expected}
\end{equation}
with $\vect{t}_{\mathrm{TT}}$ the wavefront tip-tilt.
Since the kinematic compensation has removed the rigid telescope motion measured by the encoders,
Eq.~(\ref{eq:src}) is the only path by which the mount tracking error reaches the optics.

Magnitudes along the trajectory: the GIR tracking error peaks at about $6.9\arcsec$ during the
mount start-up transient, i.e.\ $\Delta\gamma\,\theta_0 \lesssim 0.008\mas$ at $0.25\arcsec$.
Without GIR tracking the field rotates by $0.155^\circ$ in 30\,s and $1.55^\circ$ in 300\,s,
moving the off-axis image by $0.68\mas$ and $6.8\mas$, respectively.

%======================================================================
\section{Verifications}
\label{sec:verification}

Table~\ref{tab:verif} summarizes the verifications; they are detailed below. Unit tests were re-run
on 10/20/2026 for this note; the integrated model results are from the runs whose data and figures
are stored in the demo \code{data/} and \code{figures/} folders.

\begin{table}[htbp]
  \centering
  \caption{Summary of the verifications.}
  \label{tab:verif}
  \begin{tabularx}{\linewidth}{lXl}
    \toprule
    Item & Verification & Result \\
    \midrule
    Trajectory      & regeneration with \code{atp\_trajectory.py} identical to \code{atp\_trajectory.csv} & pass \\
    Tracking        & encoders vs setpoint over 300\,s (Fig.~\ref{fig:encoders}) & pass \\
    Feed-forward    & 300\,s run with/without feed-forward; axes cross-coupling (Table~\ref{tab:coupling}) & saturation removed; coupling $\le 1.8\times10^{-6}$ \\
    Kinematics      & compensated RBMs (Fig.~\ref{fig:rbm}), segment tip-tilt drift (Fig.~\ref{fig:stt}) & drift removed \\
    GIR rotation    & unit tests \code{zero\_angle}, \code{center\_segment}, \code{no\_rotation} and the deformation rotation test & see \S\ref{sec:ver-unit} \\
    Pointing        & unit tests \code{azimuth\_elevation}, \code{elevation}, \code{gir} & see \S\ref{sec:ver-unit} \\
    Optical model   & unit tests \code{pointing\_error}, \code{pointing\_rotation\_error} & pass \\
    End-to-end      & image centroid vs Eq.~(\ref{eq:expected}) (Figs.~\ref{fig:tt}, \ref{fig:tt-off}) & $\le 0.03\mas$ rms after 2\,s \\
    GIR error       & GIR tracking error exaggerated by $10^{4}$ (Fig.~\ref{fig:tt-gain}) & fitted gain 0.999 \\
    \bottomrule
  \end{tabularx}
\end{table}

\subsection{Trajectory and mount tracking}
\label{sec:ver-traj}

Running \code{scripts/atp\_trajectory.py} again produced a file identical to the committed
\code{atp\_trajectory.csv}; the script reports the target ICRS coordinates, the azimuth, elevation and
rotator ranges of Section~\ref{sec:astropy}, and an initial parallactic angle of $-180^\circ$, as
expected for a target on the meridian ($H=0$) north of a southern site.

Figure~\ref{fig:encoders} compares the 1\,Hz trajectory to the averaged mount encoders over the full
300\,s at 1\,kHz. The encoders follow the interpolated setpoints; the tracking error is a start-up
transient of the first 2\,s (up to $1.4\arcsec$ in azimuth, $0.02\arcsec$ in elevation and
$6.9\arcsec$ for the GIR) after which it is negligible, with a small GIR ripple.

\begin{figure}[htbp]
  \centering
  \includegraphics[width=\linewidth]{trajectory_vs_encoders.png}
  \caption{ATP trajectory vs mount encoders (left) and tracking error (right), 300\,s at 1\,kHz.}
  \label{fig:encoders}
\end{figure}

\subsection{Mount feed-forward}
\label{sec:ver-ff}

Without feed-forward (\code{MOUNT\_FEEDFORWARD=0}) the azimuth drives saturate at about 242\,s and the
mount stops following the trajectory. With the feed-forward of Eq.~(\ref{eq:stiffness}) the full 300\,s
trajectory runs without saturation and with the tracking error of Figure~\ref{fig:encoders}: the
integrator of the mount controller no longer has to build up the spring torque.

The cross-coupling neglected by the feed-forward is estimated from the same DC gain as
Eq.~(\ref{eq:stiffness}). The coupling matrix
\begin{equation}
  C_{k,l} = \kappa_l \sum_{j\in\mathcal{D}_l} g_{k,j}
\end{equation}
is the averaged encoder angle of axis $k$ per radian of feed-forward on axis $l$; its diagonal is 1 by
construction. For the 2025 FEM (Table~\ref{tab:coupling}) the off-diagonal terms are at most
$1.8\times10^{-6}$, and the per drive gains $g_{k,j}$ across axes are $10^{-5}$ to $10^{-7}$ of those within
an axis. The largest terms are the elevation and the GIR responses to the azimuth feed-forward, the
azimuth having the largest setpoint ($-1.70^\circ$ at 300\,s). The static errors left by the uncoupled
feed-forward at the end of the trajectory, $\sum_{l\neq k} C_{k,l}\theta_l$, are
$5.0\times10^{-8}$\,rad ($10\mas$) in elevation, $7.3\times10^{-9}$\,rad ($1.5\mas$) for the GIR and below
$1\mas$ in azimuth. They build up slowly (about $34\,\mu$as/s in elevation) and are removed by the mount
feedback for a negligible torque, about 0.5\,N\,m per elevation drive against the 800\,kN\,m saturation,
which is why no residual appears in the tracking error of Figure~\ref{fig:encoders}.
A decoupled feed-forward, $C^{-1}$, would change the torques by at most $1.8\times10^{-6}$ of the diagonal
terms and is not needed.

\begin{table}[htbp]
  \centering
  \caption{Mount axes cross-coupling $C_{k,l}$ of the feed-forward: averaged encoder angle of axis $k$
    (rows) per radian of feed-forward on axis $l$ (columns), 2025 FEM.}
  \label{tab:coupling}
  \begin{tabular}{lccc}
    \toprule
         & az                    & el                    & GIR \\
    \midrule
    az   & 1                     & $-1.8\times10^{-6}$   & $-2.4\times10^{-9}$ \\
    el   & $-1.7\times10^{-6}$   & 1                     & $-3.6\times10^{-9}$ \\
    GIR  & $-2.5\times10^{-7}$   & $4.4\times10^{-9}$    & 1 \\
    \bottomrule
  \end{tabular}
\end{table}

\subsection{Kinematic compensation}
\label{sec:ver-kin}

The kinematic compensation was reached in three steps, each verified on the segment tip-tilt
from the LOM, the most sensitive indicator of a residual rigid motion:
\begin{enumerate}
\item \emph{Geometric removal} (\code{scripts/rigid\_body.py}): the rigid motion is modeled as the
  linearized rotation of the segment vertices and frames about the azimuth axis
  $(0, \sin 30^\circ, \cos 30^\circ)$ and the elevation axis $(-1, 0, 0)$ of the 30$^\circ$ zenith FEM.
  The mismatch between this geometric model and the FEM kinematics left a segment tip-tilt drift that
  grows with the azimuth angle.
\item \emph{FEM DC gain} (Eq.~\ref{eq:kgain}) applied to the M1 and M2 rigid body motions only:
  the geometric mismatch is gone, but the segments still drift radially by up to 15--20\,mas in 50\,s
  (Figure~\ref{fig:stt}, left). The same drift is obtained with a FEM and M1 only model
  (\code{trajectory-mount} demo), which pointed to the M1 hardpoint forces: their kinematic component
  is offloaded by the M1 actuators force loop. A closed-loop calibration of the azimuth column of $K$
  was tried as a workaround.
\item \emph{Compensation of all the FEM outputs} but the mount encoders, including the control
  feedback (Section~\ref{sec:kin-feedback}): the drift disappears without any calibration. Over
  30--300\,s the segment tip-tilt stays within $\pm 1.5\mas$ with no secular trend
  (Figure~\ref{fig:stt}, right).
\end{enumerate}
With the final compensation, the M1 and M2 rigid body motions show the start-up transient
(tens of \textmu m, sub-arcsecond) decaying to zero within about 20\,s and no motion following the
1.7$^\circ$ azimuth rotation (Figure~\ref{fig:rbm}).

\begin{figure}[htbp]
  \centering
  \includegraphics[width=0.49\linewidth]{segment_tiptilt_orbits_60s_nocal.png}
  \includegraphics[width=0.49\linewidth]{segment_tiptilt_orbits_300s.png}
  \caption{Segment tip-tilt orbits. Left: FEM DC gain kinematics on the M1 and M2 rigid body motions
    only, 10--60\,s. Right: kinematic compensation of all the FEM outputs, 30--300\,s.}
  \label{fig:stt}
\end{figure}

\begin{figure}[htbp]
  \centering
  \includegraphics[width=0.9\linewidth]{rigid_body_motions.png}
  \caption{M1 and M2 segment rigid body motions with kinematic compensation along the 300\,s trajectory.}
  \label{fig:rbm}
\end{figure}

\subsection{Unit tests}
\label{sec:ver-unit}

\paragraph{\code{gmt\_dos-clients\_crseo} (\code{optical\_model/read/pointing.rs}).}
The tests use a V band camera ($\lambda/(4D)$ sampling binned $2\times2$) in absolute pointing along
the optical axis; results of the re-run:
\begin{itemize}
\item \code{pointing\_error}: on-axis source, pointing errors of $20\mas$ along $x$ then $y$. The frame
  centroid moves to $(-19.758, 0.001)\mas$ and $(0.003, -19.752)\mas$, i.e.\ by $-\vect{p}$ (the
  centroid of the truncated frame is $\sim$1\% short), and the wavefront tip-tilt stays at
  $0.000\mas$. The test passes within $5\%\,|\vect{p}| + 0.01\mas$ on the image and $0.01\mas$ on the
  tip-tilt.
\item \code{pointing\_rotation\_error}: source $100\mas$ off-axis along $x$. Without error the centroid
  is at $(98.313, -0.009)\mas$; with $\Delta\gamma = 0.04$\,rad it moves to $(98.230, -3.964)\mas$,
  i.e.\ by $(-0.083, -3.955)\mas$ against $((\cos\Delta\gamma - 1)\theta_0, -\sin\Delta\gamma\,\theta_0)
  = (-0.080, -3.999)\mas$: the image is rotated by $-\Delta\gamma$. Resetting the error restores the
  initial position.
\item \code{pointing\_scan} (ignored by default) scans the pointing error over $1.5$ pixel and compares
  the absolute pointing image motion to the one of an equivalent wavefront tilt. It revealed the CEO
  imaging axis convention fixed in \code{crseo-sys}~1.3.3.
\end{itemize}

\paragraph{\code{trajectory-servos-mount}.}
\begin{itemize}
\item \code{pointing::tests::azimuth\_elevation}: for a trajectory starting at $45^\circ$ elevation,
  encoders $(1, 2)$\,\textmu rad away from the setpoint give
  $\vect{p} = (10^{-6}\cos(45^\circ + 2\,\mu\mathrm{rad}), 2\times 10^{-6})$; a quarter turn of
  the GIR maps $x\to y$ and $y\to -x$.
\item \code{pointing::tests::elevation}: for a trajectory starting at $45^\circ$ elevation with the
  telescope $15^\circ$ lower and a 1\,\textmu rad azimuth error, $p_x = 10^{-6}\cos 30^\circ$: the
  $\cos(\mathrm{el})$ factor follows the telescope elevation.
\item \code{pointing::tests::gir}: with the GIR tracking a field rotation of 3\,\textmu rad and an encoder
  at 1\,\textmu rad, $\Delta\gamma = -2$\,\textmu rad; with the GIR at rest, $\Delta\gamma = -q$.
\item \code{gir\_rotation::tests::zero\_angle}: no rigid body motion for $\gamma = 0$, for M1 and M2.
\item \code{gir\_rotation::tests::center\_segment}: for $\gamma = 10^{-2}$\,rad, M1 S7 (on axis) has no
  translation and a rotation of exactly $\gamma$ about $z$.
\item \code{gir\_rotation::tests::no\_rotation}: for $\gamma = 0$, any deformation of the M1 and M2
  segments is returned unchanged.
\item \code{gir\_rotation::tests::deformation\_rotates\_with\_the\_mirror}: M1 S7 translated by
  1\,\textmu m along $x$ and tilted by 2\,\textmu rad about $x$, rotated by $\gamma = 90^\circ$, is translated
  along $y$ and its normal, initially tilted towards $-y$, is tilted towards $+x$.
\end{itemize}

\subsection{End-to-end verification of the image motion}
\label{sec:ver-e2e}

\code{scripts/plot\_tiptilt\_vs\_tracking.py} compares the centroid of the image on the camera
to the prediction of Eq.~(\ref{eq:expected}) built from the logged \code{PointingError},
\code{PointingRotationError} and \code{TipTilt}. Frames where the predicted image is outside 80\% of the
FOV are discarded (the centroid is biased near the frame edges) and a static centroid bias,
estimated after 2\,s, is removed. The off-axis angle $\theta_0$ of the run is given to the script
as its 6th argument (default $0.25\arcsec$).

Figure~\ref{fig:tt} is a 30\,s run with GIR tracking and an on-axis source (\code{OFF\_AXIS=0}),
so it verifies the pointing error and the structure tip-tilt terms of Eq.~(\ref{eq:expected}), the
rotation term being zero on axis. The static centroid bias is $(0.97, 0.98)\mas$, a quarter of a pixel.
During the first second the pointing error reaches about $1\arcsec$ in $x$ and the image leaves the
field. The image then follows $-\vect{p}$, and after a few seconds the structure tip-tilt dominates
over the pointing error. Re-running the script on the stored data gives the rms values of
Table~\ref{tab:e2e}: the image motion is explained to $0.27\mas$ rms in $x$ over the whole run,
the largest residuals (up to $4\mas$) occurring during the first 1.5\,s when the image is close to the
edge of the field, to $0.03\mas$ rms after 2\,s and to $0.007\mas$ after 5\,s. After 10\,s the image
motion settles to $0.1$--$0.3\mas$, dominated by the structure tip-tilt, the on-sky tracking error
being below $0.1\mas$.

Figure~\ref{fig:tt-off} is a 30\,s run with \code{GIR\_TRACKING=0}: the field rotation term grows
linearly and reaches $0.69\mas$ at 30\,s, against $0.676\mas$ expected from
$\theta_0\,|q(30\,\mathrm{s})|$, and the image motion follows it once the start-up transient
has decayed.

With GIR tracking, the rotation term is the GIR tracking error times the off-axis angle: at most
$3.3\times10^{-5}\,\mathrm{rad}\times 0.25\arcsec \simeq 0.008\mas$ during the start-up transient,
when the residual reaches several mas, and about $10^{-4}\mas$ afterwards, below the $0.002$--$0.007\mas$
residual. It cannot be measured as such. To verify the rotation path in that configuration, a 30\,s
run with GIR tracking and the source $0.25\arcsec$ off-axis was made with \code{PointingRotationError}
exaggerated by a gain of $10^{4}$ (\code{POINTING\_ROTATION\_GAIN=1e4}, Figure~\ref{fig:tt-gain}): the
rotation error peaks at 0.33\,rad, moving the image by up to $82\mas$, and the GIR tracking ripple after
the transient moves it by about $0.1\mas$. The image motion $y$ minus $(-\vect{p} + \vect{t}_{\mathrm{TT}})$
is fitted by the rotation term times a gain (Table~\ref{tab:gain}): the fitted gain is 0.999 after
5\,s, and the rotation term reduces the residual from $0.106\mas$ to $0.0013\mas$ rms. The only large
residual, 50\,mas at 0.3\,s in $x$, occurs when the pointing transient brings the off-axis image to the
edge of the field. The rotation by the GIR tracking error is therefore applied as predicted, and, the term
being linear in $\Delta\gamma$, this holds for the actual, $10^{4}$ times smaller, error.

\begin{table}[htbp]
  \centering
  \caption{Image motion along $y$ explained by the rotation term, 30\,s run with GIR tracking, source
    $0.25\arcsec$ off-axis and \code{PointingRotationError} exaggerated by $10^{4}$ (1468 out of 1500 frames
    in the field). Residual rms [mas] without and with the rotation term, and fitted gain of the term.}
  \label{tab:gain}
  \begin{tabular}{lcccc}
    \toprule
                & rotation term & residual without & residual with & fitted gain \\
    \midrule
    $t > 0$\,s  & 5.904 & 5.802 & 0.231  & 0.982 \\
    $t > 2$\,s  & 0.120 & 0.120 & 0.010  & 0.992 \\
    $t > 5$\,s  & 0.106 & 0.106 & 0.0013 & 0.9992 \\
    \bottomrule
  \end{tabular}
\end{table}

\begin{table}[htbp]
  \centering
  \caption{Image motion and residual rms [mas] of the 30\,s on-axis run with GIR tracking (1472 out of 1500
    frames in the field).}
  \label{tab:e2e}
  \begin{tabular}{lcccccc}
    \toprule
                & \multicolumn{2}{c}{image motion} & \multicolumn{2}{c}{$-\vect{p}$} & \multicolumn{2}{c}{residual} \\
                & $x$ & $y$ & $x$ & $y$ & $x$ & $y$ \\
    \midrule
    $t > 0$\,s  & 28.35 & 4.24 & 27.52 & 1.36 & 0.274 & 0.034 \\
    $t > 2$\,s  & 3.35  & 1.80 & 2.37  & 0.21 & 0.028 & 0.014 \\
    $t > 5$\,s  & 0.49  & 0.34 & 0.11  & 0.07 & 0.007 & 0.002 \\
    \bottomrule
  \end{tabular}
\end{table}

\begin{figure}[htbp]
  \centering
  \includegraphics[width=0.85\linewidth]{tiptilt_vs_tracking_onaxis.png}
  \caption{Mount tracking error and image motion, 30\,s on-axis run with GIR tracking. From top to
    bottom: tracking error, image centroid vs $-\vect{p}$, magnitudes of each term of
    Eq.~(\ref{eq:expected}), residual.}
  \label{fig:tt}
\end{figure}

\begin{figure}[htbp]
  \centering
  \includegraphics[width=0.85\linewidth]{tiptilt_vs_tracking_gir_off.png}
  \caption{Same as Figure~\ref{fig:tt} for a 30\,s run with \code{GIR\_TRACKING=0} and a source
    $0.25\arcsec$ off-axis: the image rotates with the field.}
  \label{fig:tt-off}
\end{figure}

\begin{figure}[htbp]
  \centering
  \includegraphics[width=0.85\linewidth]{tiptilt_vs_tracking_gir_gain1e4.png}
  \caption{Same as Figure~\ref{fig:tt} for a 30\,s run with GIR tracking, a source $0.25\arcsec$
    off-axis and the GIR tracking error exaggerated by $10^{4}$ (\code{POINTING\_ROTATION\_GAIN=1e4}).}
  \label{fig:tt-gain}
\end{figure}

\end{document}
